\scanpage{78}
% lines 1–5, extensions and morphisms
Quand on interprète $P,P'$ en termes d'ext. de $\underline O$ par des Modules $L,L'$, munies de $T,T'$ \unsure{telle|telles} que $\psi(T)=\unsure{\chi|\gamma}$, $\psi'(T')=\chi$, la condition de compatibilité avec les $\chi$-trivialisations s'écrit simplement
\[
f(T)=T'.
\]
% lines 6–11, chosen splitting
\uline{Cas scindé} Un scindage d'un torseur $\chi$-trivialisé $(P,\deleted{\unsure{T}}\pi)$ est, par déf., une section (i.e. trivialisation) \inserted{$x_0$} de $P$. Alors \deleted{\unsure{à}}
\[
P\xleftarrow{\sim} L,\qquad u\longmapsto x_0+u
\]
\deleted{et la donnée de la $\pi$-tri} d'où
\[
E\xleftarrow{\sim}L+\underline O,\qquad(u,\lambda)\longmapsto\lambda x_0+u
\]
% lines 12–18, base point and coordinates
La donnée de la $\chi$-trivialisation équivaut à celle de l'élément
\[
\underset{\text{ou }b_{x_0}}{b}=\pi(x_0)=\chi x_0-T
\]
de \deleted{\unsure{terme}} sorte qu'en termes de $b$ on aura
\[
T=\chi x_0-b
\]
ou en termes de l'iso $E\simeq L+\underline O$,
\[
T=(-b,\chi)
\]
% lines 19–23, linear formula and check
On aura donc, en termes de $b$
\[
\pi(\overbrace{\underbrace{\lambda x_0+u}_{(u,\lambda)}}^{x})=\lambda b+\chi u
\]
[Vérifions que c'est \unsure{aussi|vrai|OK} \deleted{$\unsure{\pi}$} $\chi x-\psi(x)T=\chi\lambda x_0+\chi u-\lambda\underbrace{T}_{\chi x_0-b}$ \deleted{\unsure{$=$}}, on trouve bien $\chi u+\lambda b$ OK ]
% archive sheet number: 77
